% ═══════════════════════════════════════════════════════════════════════════════
%  CAPÍTULO — Builder
% ═══════════════════════════════════════════════════════════════════════════════

\chapter{Builder}

\patroninfo{Creacional}{97}

% ─── Situación Problema ──────────────────────────────────────────────────────
\section{Situación Problema}

\section{Situación Problema}

El motor de búsqueda de Airbnb permite al usuario aplicar una gran
cantidad de filtros antes de explorar propiedades: destino, rango de
fechas, número de huéspedes, precio mínimo y máximo, tipo de alojamiento,
amenidades requeridas, política de cancelación y puntuación mínima, entre
otros. No todos los filtros son obligatorios; de hecho, la mayoría de las
búsquedas utilizan solo un subconjunto de ellos.

Intentar representar todas las combinaciones posibles mediante constructores
sobrecargados resulta inviable: con diez parámetros opcionales se
producen decenas de firmas distintas, muchas de ellas ambiguas cuando
varios parámetros son del mismo tipo. Pasar \texttt{null} para los campos
no utilizados es una práctica propensa a errores y difícil de leer.
Además, la consulta resultante debe poder validarse como un todo coherente
antes de ejecutarse.

% ─── Solución ────────────────────────────────────────────────────────────────
\section{Solución}

El patrón Builder separa la construcción del objeto \texttt{ConsultaBusqueda}
de su representación final. Un objeto \texttt{BusquedaBuilder} expone
métodos encadenables —\texttt{conDestino()}, \texttt{conFechas()},
\texttt{conPrecioMaximo()}, \texttt{conAmenidades()}— que van acumulando
los criterios seleccionados por el usuario. Al final, el método
\texttt{construir()} valida que los criterios sean coherentes y devuelve
un objeto \texttt{ConsultaBusqueda} completamente configurado.

Esto elimina los constructores con largas listas de parámetros nulos,
hace el código de construcción legible y autodescriptivo, y centraliza
las reglas de validación en un único lugar. Si en el futuro se añaden
nuevos filtros, solo es necesario agregar el método correspondiente al
builder sin afectar al resto del sistema.

% ─── Diagrama de clases ─────────────────────────────────────────────────────
\begin{figure}[H]
    \centering
    % \includegraphics[width=\textwidth]{imagenes/abstract_factory_diagrama.png}
    \caption{Diagrama de clases del patrón Abstract Factory}
\end{figure}


% ─── Roles de las Clases ────────────────────────────────────────────────────
\section{Roles de las Clases}

% Explicar la responsabilidad de cada clase/interfaz

\begin{itemize}
    \item \textbf{AbstractFactory:}
    \item \textbf{ConcreteFactory:}
    \item \textbf{AbstractProduct:}
    \item \textbf{ConcreteProduct:}
    \item \textbf{Client:}
\end{itemize}


% ─── Main ───────────────────────────────────────────────────────────────────
\section{Main}

% Explicar qué realiza el programa principal

\begin{lstlisting}[language=Java, caption=NombreClase]
 
\end{lstlisting}


% ─── Salida del Main ────────────────────────────────────────────────────────
\section{Salida del Main}

\begin{lstlisting}[language=Java, caption=NombreClase]
 
\end{lstlisting}


% ─── Clases ─────────────────────────────────────────────────────────────────
\section{Clases}

% ─── Clase 1 ────────────────────────────────────────────────────────────────
\subsection{NombreClase}

\begin{lstlisting}[language=Java, caption=NombreClase]
 
\end{lstlisting}


% ─── Clase 2 ────────────────────────────────────────────────────────────────
\subsection{NombreClase}

\begin{lstlisting}[language=Java, caption=NombreClase]
 
\end{lstlisting}


% ─── Clase 3 ────────────────────────────────────────────────────────────────
\subsection{NombreClase}

\begin{lstlisting}[language=Java, caption=NombreClase]
 
\end{lstlisting}


% ─── Conclusiones ───────────────────────────────────────────────────────────
\section{Conclusiones}

% Explicar:
% - Ventajas del patrón
% - Cuándo usarlo
% - Beneficios obtenidos
% - Posibles desventajas
% - Relación con principios SOLID